\chapter{Photometric (Camera) Systems}
\label{ch:camera}

\section{Imaging Architectures}
\label{sec:camarch}

Photometric launch monitors estimate motion from calibrated images taken near impact.
Their useful capture volume depends on camera placement, illumination, exposure timing, resolution and visibility; it is not a universal length of ball flight.
Multiple cameras can supply simultaneous geometric constraints, while strobed exposures can supply several time-separated observations within one image.
Neither camera count nor a nominal frame rate alone specifies measurement uncertainty.

Commercial examples illustrate different observation geometries.
Foresight's GC3 uses three cameras, while the GCQuad family uses four; supported club outputs depend on the model and enabled features~\cite{foresightgc3,foresightgcquad}.
Uneekor's overhead EYE~XO2 uses three cameras and club stickers for the documented club-face data; overhead placement does not itself make club measurement markerless~\cite{uneekorxo2manual}.
ProTee VX instead advertises two overhead cameras and sticker-free club and ball outputs~\cite{proteevx}.
Garmin's R50 uses three cameras and a club-face sticker for club tracking~\cite{garminr50,garminr50support,garminr50stickers}.
These are manufacturer descriptions of capabilities, not a common independent accuracy specification.
The model-specific distinctions in \cref{ch:survey} matter more than a camera-versus-radar label, especially for hybrid instruments.

\section{Ball Measurement}

\subsection{Photogrammetric Position and Velocity}

A calibrated camera maps a 3D point to an image through its intrinsics, pose and lens-distortion model.
For matched observations of the \emph{same physical point}, corrected image coordinates can be triangulated from several views, for example by minimizing reprojection error.
A ball silhouette needs additional care: its apparent contour is not a set of fixed material points shared between viewpoints.
Estimating the sphere center from its outline should therefore use a sphere-projection model before treating the result as a point observation.

Let $\vect{p}_k$ and $\vect{p}_{k+1}$ be estimated centers in one calibrated target frame at times separated by $\Delta t>0$.
Then
\begin{equation}
\overline{\vect{v}}_k=\frac{\vect{p}_{k+1}-\vect{p}_k}{\Delta t}
\end{equation}
is an interval velocity estimate; its norm is speed.
Displacement alone is not speed, and an interval estimate need not equal the instantaneous velocity at separation from the club.
A local trajectory fit can estimate velocity at a declared event time, with timing uncertainty and any extrapolation included in the error budget.
Launch angles follow by resolving that velocity in the target frame, not by reading image-axis slopes as physical launch angles.

For an ideal rectified stereo pair, let $B$ be the baseline, $f$ the common focal length in pixels, and $d=u_L-u_R>0$ the disparity in pixels.
Depth along the camera optical axis is $Z=fB/d$.
Holding calibration fixed, first-order propagation gives
\begin{equation}
\label{eq:depth}
\sigma_Z\simeq\frac{Z^2}{fB}\sigma_d,
\qquad
\sigma_d^2=\operatorname{Var}(u_L)+\operatorname{Var}(u_R)
-2\operatorname{Cov}(u_L,u_R).
\end{equation}
Thus two independent image-coordinate errors of standard deviation $\sigma_{px}$ give $\sigma_d=\sqrt{2}\sigma_{px}$, not $\sigma_{px}$.
The approximation becomes unreliable when disparity uncertainty is large relative to disparity.
Baseline, focal-length, distortion and alignment uncertainty add further contributions; \cref{eq:depth} is not a bound for every multi-camera layout.
Camera depth $Z$ must also be transformed into the book's target frame before comparison with launch-monitor coordinates.

\subsection{Why an Ellipse Center Is Not the Ball Center}

The distinction can be derived directly for an ideal pinhole camera.
Let a sphere of calibrated radius $r$ have center $\vect{c}=(X,Y,Z)^\top$ in camera coordinates, with $Z>r$.
For a normalized image ray $\vect{q}=(u,v,1)^\top$, tangency requires
\begin{equation}
(\vect{c}^{\top}\vect{q})^2
=(\vect{c}^{\top}\vect{c}-r^2)\vect{q}^{\top}\vect{q}.
\end{equation}
The center of this image conic is
\begin{equation}
(u_e,v_e)=\frac{Z}{Z^2-r^2}(X,Y),
\qquad
(u_c,v_c)=\frac{1}{Z}(X,Y)
\end{equation}
for the ellipse center and projected sphere center, respectively.
Their difference is $r^2/(Z^2-r^2)$ times $(u_c,v_c)$.
It vanishes on the optical axis but generally does not vanish merely because an ellipse was fitted accurately.
These are normalized coordinates; the intrinsics convert them to pixels.
A sphere model, calibration and known or estimated radius supply the correction, while partial occlusion, ball deformation near contact and nonuniform illumination require separate treatment.
The familiar ball-size range cue is therefore conditional on size and projection assumptions, not an independent exact depth measurement.

\subsection{Spin: Surface-Pattern Registration}
\label{sec:dimplespin}

Visible dimples, logos or printed marks can constrain rotation by matching their surface locations between exposures.
A geometric estimator jointly needs the ball center, radius, camera calibration, exposure times and feature correspondences.
Changes in lighting, specular reflections, occlusion and repeated texture can produce matches that are photometrically plausible but geometrically wrong.
Describing this as spherical registration explains a possible observation model; it does not establish the proprietary algorithm inside every camera product.

Let $Q_k\in SO(3)$ map ball-fixed coordinates into a fixed world frame.
The spatial rotation increment is $R=Q_{k+1}Q_k^\top$.
Under constant spatial angular velocity over $\Delta t$, and on the branch $\|\vect{\omega}\|\Delta t<\pi$,
\begin{equation}
\label{eq:rotangle}
\vect{\omega}=\frac{\operatorname{vee}(\operatorname{Log}R)}{\Delta t},
\qquad
\theta=\arccos\!\left(\frac{\operatorname{tr}R-1}{2}\right)\in[0,\pi].
\end{equation}
Here $\operatorname{vee}$ converts a skew matrix to its axial vector and $\operatorname{Log}$ denotes the principal rotation logarithm on that branch.
The trace angle alone is unsigned; an eigenvector with eigenvalue one also has an arbitrary sign.
A consistent rotation-vector construction supplies the orientation information away from the branch boundary.
Near zero rotation the axis direction is poorly determined; near $\pi$ the principal branch needs particular care.
If angular velocity changes direction during the interval, the logarithm describes an equivalent finite rotation and is not generally its arithmetic time average.

Rotations separated by complete revolutions can generate the same endpoint orientation, and repeated surface patterns can introduce further ambiguity.
Additional exposures, independent observations or a justified speed prior may select a candidate; they do not guarantee that an otherwise ambiguous rotation becomes observable.
Three noncollinear matched 3D surface points are sufficient for a noiseless proper rigid-transform fit, but three arbitrary image detections are not automatically three known 3D correspondences.
A marked ball can improve correspondence without removing timing or calibration error.
Rapsodo's RPT pattern is an optical spin aid; Titleist RCT's embedded radar-reflective structure serves a different observation mechanism~\cite{rapsodofaq,garminr10rct}.
No universal percentage spin accuracy, or independence from ball speed, follows from the registration equations.

\section{Club Measurement}
\label{sec:camclub}

\subsection{Fiducial Geometry and Reference Points}

Reflective club markers provide identifiable image features.
Foresight's documented single-marker mode supplies a limited club-data set, while supported four-marker modes add orientation and impact information~\cite{foresightmarkers,foresightquadmanual}.
The actual marker layout matters: four collinear points still leave rotation about their line unconstrained.
A suitable noncollinear constellation can define a rigid pose when its geometry and correspondences are known and visible.
For a curved driver face, the relationship between that constellation and the selected face-center tangent plane also needs calibration; a plane fitted to arbitrary surface points is not automatically the desired face plane.

A single marker trajectory provides that point's position and velocity, not a unique full club pose.
For two points $P$ and $O$ fixed in a rigid clubhead,
\begin{equation}
\vect{v}_P=\vect{v}_O+\vect{\omega}\times(\vect{p}_P-\vect{p}_O),
\end{equation}
with every quantity expressed in the same frame and at the same time.
Consequently, differentiating a pose origin or marker position does not automatically yield the manufacturer's club-speed reference point.
Path and attack angle inherit this point and timing dependence.
Impact coordinates additionally require a registered face frame and a defined ball-contact event; a replay image alone is not a calibrated numerical impact-location measurement.
These are the same measurand distinctions that govern the radar comparison in \cref{ch:accuracy}.

The GC3's published club list contains speed, path, attack angle and smash factor; it should not be credited with the GCQuad's face-angle and impact-location feature set~\cite{foresightgc3}.
Nor is a ratio such as smash factor a separate directly observed primitive simply because it appears in a measured-data list.
Feature entitlement, marker mode and metric definition should accompany any comparison between units.

\subsection{Overhead Tracking and Calibration}

Overhead placement offers useful views but does not guarantee an unobstructed face, a marker-free algorithm or a uniquely determined pose.
The EYE~XO2 manual specifies club stickers, whereas ProTee VX advertises sticker-free image-based club outputs~\cite{uneekorxo2manual,proteevx}.
Those different designs should not be collapsed into one generic overhead measurement method.
Learned shape models can help interpret images; their training domain and failure behavior remain part of the measurement system.

Factory calibration concerns internal camera geometry; setup also establishes the target line, ground plane, working distance and hitting zone.
Moving or tilting a unit, changing its support or obscuring calibration features can change that relationship.
The EYE~XO2, for example, specifies a mounting envelope and a calibration procedure~\cite{uneekorxo2manual}; other products require their own instructions.
A DIY build should verify residuals on independent calibration poses and assess timing, distortion, blur and occlusion across its intended volume.
Fitting the calibration images well is not a held-out accuracy test.

\section{The Indoor Argument and Downrange Limits}

A camera system that obtains the necessary launch and spin observations near impact can estimate launch conditions before the ball reaches a screen.
This is an important indoor advantage, conditional on the capture window and visibility actually being adequate.
It does not observe the subsequent aerodynamic trajectory, environmental variation, landing or bounce beyond that screen.
Those outputs are predictions from the estimated state and the flight/ground models of \cref{ch:flight}.
Short-flight radar systems likewise combine observations with estimation, while hybrid systems can add complementary image information.
The useful comparison is which quantities are identifiable under the installed geometry and which uncertainty remains, rather than a claim that one modality measures everything the other models.

\section{DIY and Open-Source Photometric Systems}
\label{sec:diy}

PiTrac documents Raspberry Pi cameras, infrared strobes and image processing for launch speed, angles and three-axis spin~\cite{pitrac}.
Strobed multi-exposure capture can separate effective observation times from the camera's frame rate: pulse duration controls motion blur, while pulse spacing controls displacement, overlap and rotational ambiguity.
Its evolving build documentation should be consulted for the actual camera configuration, calibration and software version; a computer-and-camera subtotal is not a complete installed-system cost.
Neither a parts list nor agreement between two frames establishes a validated error distribution across balls, speeds and lighting conditions.
The shared imaging principles discussed in \cref{sec:wintriss} establish neither direct product lineage nor rights clearance.

\begin{implication}
Optical launch and club measurements depend on a connected chain: visible features, calibrated geometry, synchronized exposures, defined reference points and justified estimation assumptions.
A change anywhere in that chain can look like a change in the golfer's swing.
Preserving those definitions and validating across the intended operating conditions makes the resulting motion estimates interpretable; adding a camera alone does not establish accuracy or all club-face capabilities.
\end{implication}
